We now fold the deposited-energy spectra of Secs.~\ref{sec:cevns} and~\ref{sec:backgrounds} through the non-paralyzable response matrix $R(E_{\rm rec}\,|\,E_{\rm dep})$ of Sec.~\ref{sec:response} to obtain the reconstructed-energy differential rate spectra $dR/dE_{\rm rec}$ for the two trapping designs, in the assumed-shield scenario of Sec.~\ref{sec:backgrounds}: the external muon, gamma, and neutron channels enter suppressed at the NUCLEUS-demonstrated performance level, and the Ge-intrinsic cosmogenic floor enters unsuppressed. This is the payoff of the forward model: the reconstructed axis is what a bandwidth-limited readout actually records, and it differs qualitatively from the deposited axis wherever the tunneling burst saturates the 25~kHz ceiling. The fold is a discrete matrix product,
\begin{equation}
\begin{aligned}
\frac{dR}{dE_{\rm rec}}\bigg|_{j} &= \frac{1}{\Delta E_{{\rm rec},j}}\sum_{i} R_{ji}\,N_{{\rm dep},i},\\
N_{{\rm dep},i} &= \frac{dR}{dE_{\rm dep}}\bigg|_{i}\,\Delta E_{{\rm dep},i},
\end{aligned}
\label{eq:fold}
\end{equation}
where $N_{{\rm dep},i}$ is the rate in deposited bin $i$, $R_{ji}$ is the probability that a deposit in bin $i$ is reconstructed into bin $j$, and $\Delta E_{{\rm dep},i}$, $\Delta E_{{\rm rec},j}$ are the bin widths. Because every column of $R$ is normalized to unit sum, Eq.~\eqref{eq:fold} conserves counts per channel by construction; we verify this numerically to a relative $\sim\!10^{-16}$ for every channel and both designs. The CEvNS and neutron-elastic spectra, computed on recoil grids in Secs.~\ref{sec:cevns} and~\ref{sec:neutrons}, are first rebinned onto the shared $E_{\rm dep}$ grid by a piecewise log--log power-law integral that conserves counts to $9.6\times10^{-5}$ overall. Because the reconstructed energy is calibrated to estimate the deposit (Sec.~\ref{sec:response}), a sub-onset deposit is reconstructed via the linear mapping $E_{\rm rec}\approx E_{\rm dep}$ rather than discarded, and carried into the matching $E_{\rm rec}$ bins.

Figure~\ref{fig:spectra} is the central result: $dR/dE_{\rm rec}$ versus reconstructed energy for the reactor-CEvNS signal and every background channel---ambient neutron elastic, cosmic muons, environmental gammas, and the Ge-intrinsic $^{3}$H/$^{68}$Ge/$^{65}$Zn floor---with the external channels shown at their assumed-shield suppression. The reactor-CEvNS signal is drawn unsuppressed; the grey band marks the reconstructed saturation region. Four features organize the spectrum.

First, and decisively, the reactor-CEvNS signal does not have the tens-of-eV region to itself. The signal, whose recoils sit entirely below the saturation onset, reconstructs to tens of eV: its differential rate has a median near $55$~eV in $E_{\rm rec}$, with about two-thirds of its counts below $100$~eV, and it integrates to $\sim$$1.19\times10^{2}$ counts~kg$^{-1}$day$^{-1}$, of which $62.6$ (Ta$\to$Al) and $63.6$ (Al$\to$Hf) fall in the $10$--$100$~eV region of interest. But the ambient neutron-elastic recoils populate exactly the same band, and even after the assumed neutron suppression factor of $200$ they integrate to $23.9$ and $24.2$ counts~kg$^{-1}$day$^{-1}$ in the RoI---the single largest background term, $\sim$99\% of the residual budget. The resulting in-band signal-to-background ratio is $S/B\simeq2.6$ for both designs, i.e.\ the signal exceeds the summed background by a factor of a few. This is the load-bearing result of the paper, and it is the reverse of what a deposited-threshold argument alone would suggest: recovering reactor CEvNS is limited not by the readout but by a nuclear-recoil background that shares the signal's kinematics. This confirms, now on the reconstructed axis of a bandwidth-limited readout, the neutron-limited character already established for reactor-CEvNS searches~\cite{biffl2023,conusneutron2019,ricochetneutron2023,nucleusbkg2026}, and the achievable $S/B$ is set by how well the ambient neutron field can be shielded and known. We attach an order-of-magnitude label to the neutron rate, so $S/B\simeq2.6$ should be read as ``of order a few,'' not as a percent-level number.

Second, the cosmic-muon deposits all reconstruct to a single saturation plateau, well clear of the signal band. Muon energy depositions span roughly $1.5$~MeV to $\sim$$197$~MeV (Sec.~\ref{sec:backgrounds}); every one drives the peak tunneling rate orders of magnitude above the $25$~kHz ceiling, so the count-integral estimator saturates and the entire channel piles up at $E_{\rm rec}\approx37.6$~keV (Ta$\to$Al) and $\approx29.9$~keV (Al$\to$Hf). We emphasize that this tens-of-keV feature is an \emph{instrumental artifact of the bandwidth ceiling, not a physical spectral line}: the deposited muon spectrum is broad and structureless, and the pile-up appears only because the saturating response compresses three-to-four orders of magnitude of deposited energy onto the bounded plateau. The saturated-regime response shape has no literature anchor at any energy and is validated by limiting cases only (Sec.~\ref{sec:response}); the muon channel lies almost entirely within this regime, so its plateau energy must be read as a modeled ceiling. Under the assumed muon veto it contributes negligibly to the RoI.

Third, the environmental-gamma Compton continuum fills the region between the two, with reconstructed peaks near $33.5$~keV (Ta$\to$Al) and $26.6$~keV (Al$\to$Hf) and an unsuppressed integrated rate of $2.1\times10^{5}$ counts~kg$^{-1}$day$^{-1}$; the incoherent-scattering binding suppression (Sec.~\ref{sec:backgrounds}) rolls the deposited continuum off near $30$--$50$~eV $E_{\rm dep}$, so no hard cut appears at the grid edge, and under the assumed lead shielding it too is negligible in the RoI.

Fourth, beneath the shielded ambient channels lies the Ge-intrinsic cosmogenic floor, which no shield attenuates. The $^{68}$Ga and $^{65}$Cu electron-capture lines image near their deposit energies: the $^{68}$Ga M line at $158.7$~eV reconstructs to $\sim$124~eV---in the shoulder just \emph{above} the RoI---while the $^{65}$Cu M line at $120$~eV reconstructs to $\sim$96~eV, just inside it, and the L and K lines of both stand at $\sim$0.75--0.87 and $\sim$4.3--4.8~keV. Together with the low-energy tail of the $^{3}$H $\beta$ continuum, the intrinsic floor contributes $\sim$0.1--0.2 counts~kg$^{-1}$day$^{-1}$ to the RoI---two orders of magnitude below the shielded neutron term today, but the irreducible limit that would remain if the ambient neutron field were shielded away entirely. It sets the ultimate floor of the scenario and can be lowered only by minimizing surface exposure, not by shielding.

The saturating muon and Compton channels dominate the raw event rate, which raises the question of whether reconstruction remains operable. The total in-wafer event rate is $1.634$~Hz ($1.366$~Hz muon $+\,0.267$~Hz Compton, before veto), giving a mean inter-event time of $0.61$~s, vastly longer than the $\sim$$40~\mu$s pulse. The resulting pileup occupancy is $R\tau\approx3.3\times10^{-5}$ at the $20~\mu$s ($50$~kHz) sampling interval ($6.5\times10^{-5}$ at the $40~\mu$s resolving time). Quiescent reconstruction is therefore \emph{not} precluded, and the occupancy remains well below unity even at ten times the assumed flux.

The two designs give qualitatively similar spectra but differ in where saturation sets in. Because the Hf trap has the larger tunneling response per unit deposit, it saturates earlier and plateaus at a lower reconstructed energy than Ta$\to$Al, shifting the saturated features (muon pile-up, Compton peak, whole-array plateau) downward in $E_{\rm rec}$; below the crossover the two coincide on the linear $E_{\rm rec}\approx E_{\rm dep}$ mapping, and their in-RoI $S/B$ agree to within the order-of-magnitude label. The shaded bands in Fig.~\ref{fig:spectra} propagate the dominant systematics through Eq.~\eqref{eq:fold}: the CEvNS $1\sigma$ reactor-flux band ($3.4\%$ above 95~eV$_{\rm nr}$, rising to $6.2\%$/$9.8\%$ at 50/20~eV$_{\rm nr}$), the order-of-magnitude neutron band, the environmental-gamma factor-of-two site-flux band, and the muon $\pm30\%$ normalization band. Of these it is the neutron band, not the readout, that dominates the uncertainty on $S/B$.
